% Member 1 - final report, III.B Dataset, Splits and the Frozen Style Pool
\subsection{Dataset, Splits and the Frozen Style Pool}
\label{sec:dataset}

All experiments use PACS~\\cite{li2017deeper}, which contains four domains
(Photo, Art Painting, Cartoon, Sketch) over seven object classes:
\texttt{dog}, \texttt{elephant}, \texttt{giraffe}, \texttt{guitar},
\texttt{horse}, \texttt{house} and \texttt{person}. The protocol is
single-source: Photo is the only labelled training domain, and the three
other domains are used exclusively for final evaluation.

\paragraph{Splits.}
We froze the partition once, in a single directory that every member reads
and no member writes to. Photo is split into 1{,}349 training, 150
validation and 171 test images; the three target domains are held out
entirely, with 2{,}048 Art Painting, 2{,}344 Cartoon and 3{,}928 Sketch
test images. The validation split is used only for reporting the validation
accuracy during training; no target-domain image and no validation image
ever influences a parameter, a checkpoint or a hyper-parameter choice. We
verified that no path appears in more than one split, and that no two
images in different splits share identical file contents or identical
decoded pixels.

\begin{table}[t]
\caption{Data splits and their experimental roles.}
\label{tab:splits}
\centering
\small
\begin{tabular}{@{}lrp{4.4cm}@{}}
\toprule
Split & Images & Purpose \\
\midrule
Photo Training    & 1,349 & Update classifier parameters; provide labeled content images for augmentation \\
Photo Validation  & 150   & Report the validation accuracy during training \\
Photo Test        & 171   & Measure in-domain accuracy \\
Art Painting Test & 2,048 & Measure cross-domain accuracy on art paintings \\
Cartoon Test      & 2,344 & Measure cross-domain accuracy on cartoons \\
Sketch Test       & 3,928 & Measure cross-domain accuracy on sketches \\
\bottomrule
\end{tabular}
\end{table}

\paragraph{Class balance.}
The seven classes are not equally represented in PACS, and the imbalance
differs between domains. We therefore report the accuracy per target domain
and additionally provide a per-class breakdown (Fig.~\ref{fig:per_class})
rather than a single pooled number, so that a gain cannot hide a collapse in
one class. The cross-domain average we report is the unweighted mean of the
three domain accuracies, which weights the three domains equally regardless
of their size.

\paragraph{The style pool.}
The style references are ten external paintings, none of which belongs to
PACS. They are frozen in a directory with a manifest, and the pool is
hashed; every generated image records the hash of the pool it was produced
with, so a run can be checked against the pool it claims to have used. The
pool is split into two disjoint half-pools of five paintings each, which
are used to test how sensitive the results are to the particular set of
styles (Section~\ref{sec:pool_sensitivity}).

The pool is external for a reason that is easy to state and easy to get
wrong. If a style reference were taken from PACS, a training image could
carry a fragment of an evaluation image, and the resulting accuracy would
not measure generalisation. Because the pool is external, no target-domain
image reaches the training set through the style channel either. The
leakage check that accompanies the pool verifies this for every style and
every target domain.

\paragraph{The content--style pairing.}
Each training image is assigned exactly one style reference. The assignment
is deterministic, spreads the ten styles as evenly as the training set
allows, and is checked to be independent of the class label: a style must
not be systematically paired with one class more often than with another,
or the classifier could learn the style as a shortcut for the class. The
fairness checker verifies this balance, together with the pool membership
and the label agreement of every generated row, before any training run is
allowed to start. A run whose manifest fails the check is refused rather
than trained.
